\documentclass[11pt]{article}
\usepackage[margin=1in]{geometry}
\usepackage{amsmath,amssymb,amsthm,mathtools}
\usepackage{booktabs}
\usepackage{enumitem}
\usepackage{microtype}
\usepackage[hidelinks]{hyperref}

\newcommand{\A}{\mathcal A}
\newcommand{\U}{\mathcal U}
\newcommand{\Q}{\mathcal Q}
\newcommand{\D}{\mathfrak D}
\newcommand{\R}{\mathbb R}
\newcommand{\dd}{\mathrm d}
\newcommand{\Lie}{\mathcal L}
\newtheorem{definition}{Definition}
\newtheorem{proposition}{Proposition}
\newtheorem{theorem}{Theorem}

\title{V22 P3--T02: A Source-Native Six-Channel Transport Candidate\\without Hidden Geometry}
\author{The \AE ther-Flow Research Project}
\date{2026-08-09}

\begin{document}
\maketitle

\begin{abstract}
This draft/control manuscript constructs one deterministic local dynamics law
for the proposal-only B1 six-scalar source state.  The law uses only a source
evolution label and six source vector fields; it uses no metric, volume form,
connection, action, target atlas, physical clock, or GR target fit.  We derive
the equations, exact local flow solution, admissible data, regular nonconstant
background, linearization, principal polynomial, conditional identities, and
a conditional P7 token update.  The construction supplies source-derived
principal structure but does not select a unique physical cone, scale, or
effective metric.  The coefficient datum remains proposal-only primitive
debt; no ontology or source law is adopted and no Distance-to-GR status is
promoted.
\end{abstract}

\section{Claim boundary and inherited state}

Let $\A$ be the inherited smooth four-dimensional primitive source arena.  It
is source-side scaffolding debt, not an already interpreted physical
spacetime.  On each open $U\subset\A$, the sealed P3--T01 B1 state is
\begin{equation}
  q=(q^1,\ldots,q^6)\in\Q_{B1}(U)=C^2(U,\R^6).
  \label{eq:B1state}
\end{equation}
The regular operational locus used below is
\begin{equation}
  q^6\ne0,
  \qquad
  \dd\rho_1\wedge\dd\rho_2\ne0,
  \qquad
  \rho_i=\frac{q^i}{q^6},
  \qquad
  p_i=\frac{1}{1+e^{-\rho_i}}quad(i=1,2).
  \label{eq:regular}
\end{equation}
The ratios and $p_i$ are conditional source observables.  They are not target
coordinates, detector readings, or realistic matter variables.

This packet chooses a deterministic differential law rather than an action.
That choice is deliberate: an action density would require a source
integration measure not provided by the inherited B1 syntax.  No measure is
silently introduced here.

\section{Proposal-only coefficient datum}

\begin{definition}[Admissible transport datum]
An admissible coefficient datum on $U$ is
\begin{equation}
  \D_U=(\tau,X_1,\ldots,X_6),
  \label{eq:datum}
\end{equation}
where $\tau\in C^3(U,\R)$ has nowhere-vanishing differential and each
$X_A\in\Gamma_{C^2}(TU)$ obeys
\begin{equation}
  \dd\tau(X_A)=1,
  \qquad A=1,\ldots,6.
  \label{eq:normalization}
\end{equation}
The label $\tau$ orders a local source initial-value problem.  It is not a
physical clock.  The $X_A$ are ordinary source tangent vectors, not a frame,
coframe, unit flow, or target causal structure.  The datum is fixed before any
comparison with GR or other target geometry.
\end{definition}

The normalization in \eqref{eq:normalization} is a parameter convention.  It
guarantees that every $X_A$ is transverse to each level set of $\tau$ and that
its local flow parameter equals the change in $\tau$.  It neither compares
vector lengths nor contracts two derivatives.  All seven objects in
$\D_U$ are new proposal-only source-extension data.  The current ontology does
not derive or adopt them.

\section{Dynamics and passive source naturality}

\begin{definition}[Six-channel source transport law]
For a fixed admissible datum $\D_U$, define
\begin{equation}
  E_A[\D_U,q]:=X_A(q^A)=0,
  \qquad A=1,\ldots,6,
  \label{eq:eom}
\end{equation}
with no sum over $A$.  Equivalently, the equation map is
$E_{\D_U}:C^2(U,\R^6)\to C^1(U,\R^6)$.
\end{definition}

Equation \eqref{eq:eom} is local and first order.  A vector field acting on a
scalar is defined without a metric or affine connection.  There is no free
scalar coefficient: all differential coefficients are the declared $X_A$.
There is no cross-channel contraction, determinant, inverse matrix, volume
form, Hodge star, divergence, or unit normal.

\begin{proposition}[Passive source-chart naturality]
Let $\varphi:U'\to U$ be a source-chart transition.  With scalar pullback
$q'^A=q^A\circ\varphi$, $\tau'=\tau\circ\varphi$, and pulled-back vector
fields $X'_A=(\dd\varphi)^{-1}X_A\circ\varphi$, one has
\begin{equation}
  X'_A(q'^A)=\bigl(X_A(q^A)\bigr)\circ\varphi,
  \qquad
  \dd\tau'(X'_A)=1.
\end{equation}
Thus the law is independent of its source-coordinate representation.
\end{proposition}

\begin{proof}
Both statements are the chain rule for a scalar and a covector paired with a
vector.  No target diffeomorphism or active physical gauge law is inferred.
\end{proof}

\section{Initial data, boundary data, and local well-posedness}

Fix $t_0$ in the image of $\tau$ and let
$\Sigma_{t_0}=\tau^{-1}(t_0)$.  By \eqref{eq:normalization}, each $X_A$ is
transverse to $\Sigma_{t_0}$.  The initial datum is one scalar
$f_A\in C^2(V_A,\R)$ for each channel on an open
$V_A\subset\Sigma_{t_0}$:
\begin{equation}
  q^A|_{V_A}=f_A.
  \label{eq:initialdata}
\end{equation}
No normal derivative is prescribed.

\begin{theorem}[Exact local characteristic solution]
Let $\D_U$ be admissible and let $\Phi_A^t$ be the local flow of $X_A$.
For sufficiently small $|t|$, the maps
$(t,y)\mapsto\Phi_A^t(y)$ are local product charts about
$V_A\subset\Sigma_{t_0}$.  The unique local $C^2$ solution of
\eqref{eq:eom}--\eqref{eq:initialdata} is
\begin{equation}
  q^A\bigl(\Phi_A^t(y)\bigr)=f_A(y).
  \label{eq:flowsolution}
\end{equation}
Restriction and compatible gluing of such solutions are componentwise, so
the fixed-datum solution assignment is a sheaf.
\end{theorem}

\begin{proof}
The $C^2$ flow theorem and transversality give the local product charts.
Along a characteristic,
$\frac{\dd}{\dd t}q^A(\Phi_A^t(y))=X_A(q^A)=0$, proving existence.
Every point in the product neighborhood lies on one characteristic meeting
$\Sigma_{t_0}$ once, proving uniqueness.  Locality of the scalar equations
gives restriction and gluing.
\end{proof}

For a source domain with boundary defined locally by a source scalar $b$, an
optional lateral datum for channel $A$ is allowed only on the inflow subset
$\{b=0,\ \dd b(X_A)<0\}$ relative to the chosen domain orientation.  No
metric normal or boundary measure is used.  Prescribing independent outflow
data generally overdetermines the characteristic solution and is rejected.

The local well-posedness claim is limited to fixed $C^2$ coefficient data and
$C^2$ componentwise initial data: existence, uniqueness, and continuous
dependence follow in the usual local flow charts.  It is source-evolution
well-posedness, not a claim of physical hyperbolicity or causality.

\section{Explicit regular background}

On a source chart with labels $(s^0,s^1,s^2,s^3)$, set $\tau=s^0$ and
\begin{align}
 X_1&=\partial_0+\partial_1,&
 X_2&=\partial_0+\partial_2,&
 X_3&=\partial_0+\partial_3,\nonumber\\
 X_4&=\partial_0+\partial_1+\partial_2,&
 X_5&=\partial_0+\partial_2+2\partial_3,&
 X_6&=\partial_0+2\partial_1+\partial_3.
 \label{eq:witnessvectors}
\end{align}
These rational source components are an exact fixture, not target-fitted
numbers.  They obey $\dd\tau(X_A)=1$ and are pairwise distinct.

For constants $a_1,a_2,c_3,c_4,c_5,c_6$ and nonzero
$\epsilon_1,\epsilon_2,c_6$, define
\begin{equation}
 \bar q^1=a_1+\epsilon_1(s^1-s^0),\quad
 \bar q^2=a_2+\epsilon_2(s^2-s^0),\quad
 \bar q^3=c_3,\quad \bar q^4=c_4,\quad
 \bar q^5=c_5,\quad \bar q^6=c_6.
 \label{eq:background}
\end{equation}
Then $X_A(\bar q^A)=0$ for every $A$, and
\begin{equation}
 \dd\bar\rho_1\wedge\dd\bar\rho_2
 =\frac{\epsilon_1\epsilon_2}{c_6^2}
  (\dd s^1-\dd s^0)\wedge(\dd s^2-\dd s^0)\ne0.
 \label{eq:rankwitness}
\end{equation}
Thus the law has a nonconstant solution on the inherited regular operational
locus.  The two response directions remain independent; the construction does
not collapse the state to one amplitude.

In the same chart, write $X_A=\partial_0+v_A^j\partial_j$.  Arbitrary local
data $f_A$ give the explicit solution
\begin{equation}
  q^A(s^0,s)=f_A\bigl(s-v_A s^0\bigr),
  \label{eq:explicit}
\end{equation}
after translating $t_0$ to zero.  Direct substitution verifies
\eqref{eq:eom}.

\section{Linearization and principal structure}

Because the coefficient datum is independent of $q$, the linearization at
any background is
\begin{equation}
  (\delta E)_A=X_A(\delta q^A).
  \label{eq:linearization}
\end{equation}
For a source covector $\xi\in T_x^*\A$, the principal symbol is
\begin{equation}
  \sigma_{\D}(\xi)_A{}^B
  =\delta_A{}^B\,\xi(X_A),
  \qquad
  \mathcal P_{\D}(\xi)
  =\det\sigma_{\D}(\xi)
  =\prod_{A=1}^{6}\xi(X_A).
  \label{eq:principal}
\end{equation}
The polynomial is homogeneous of degree six.  In the explicit fixture the six
linear factors are distinct.  This discharges the narrow P3--T02 requirement
to generate source-derived principal structure.  It also gives the decisive
limitation: a product of six unrelated linear factors does not, without a new
source-side reduction theorem or selector, determine one quadratic cone,
physical scale, or Lorentzian $g_{\mathrm{eff}}$.

Relative to the proposal-only label $\tau$, the system is diagonal with real
source characteristic velocities and is locally strongly solvable.  Calling
this physical hyperbolicity, causality, or a light cone would exceed the claim
boundary.

\section{Constraints, identities, and conditional token update}

The coefficient constraints are $\dd\tau\ne0$, $X_A\in C^2$, and
$\dd\tau(X_A)=1$.  The operational solution constraints are
\eqref{eq:regular}.  There is no algebraic field constraint, gauge constraint,
secondary constraint, or cross-channel compatibility equation.  Regularity
is an open condition and can be lost outside a chosen local patch; the law
does not automatically preserve it globally.

For every $C^1$ scalar function $h_A$, the chain rule gives the conditional
identity
\begin{equation}
  X_A\bigl(h_A(q^A)\bigr)=h_A'(q^A)X_A(q^A)=0.
  \label{eq:balance}
\end{equation}
This is an on-equation transport identity.  Without a source measure or
divergence it is not an integrated conservation law and is not an independent
physical postulate.

Since $X_6(q^6)=0$ and $X_i(q^i)=0$ for $i=1,2$, the inherited ratios obey
\begin{align}
 X_6(\rho_i)
 &=\frac{(X_6-X_i)q^i}{q^6},\nonumber\\
 X_6(p_i)
 &=p_i(1-p_i)\frac{(X_6-X_i)q^i}{q^6},
 \qquad i=1,2.
 \label{eq:tokenupdate}
\end{align}
Equation \eqref{eq:tokenupdate} is a derived conditional continuum update for
the two protected P7 token parameters.  It is not a continuum limit theorem,
matter equation, detector law, universal coupling law, or adopted physical
conservation law.

\section{Coefficient provenance, dimensions, and source-purity audit}

The coefficient datum has exactly one provenance class: new proposal-only
source-extension data introduced by this packet.  The witness components in
\eqref{eq:witnessvectors} are small rational test values chosen to make the
construction executable; they are not inferred from a target metric or fit to
GR.  The general theorem uses arbitrary admissible $X_A$, so no fixture value
is promoted to a law constant.

Let the source flow parameter and $\tau$ carry a formal label unit $T_s$, and
let $q^A$ carry a formal channel unit $Q_A$.  Then $X_A(q^A)$ has unit
$Q_A/T_s$ and $\dd\tau(X_A)$ is dimensionless.  The ratios $\rho_i$ and
$p_i$ require $Q_i=Q_6$ on the operational patch.  These are bookkeeping
units, not physical time, length, mass, or detector calibration.

The source-purity audit finds no metric, inverse metric, determinant, volume
form, connection, Hodge operator, derivative contraction, norm, unit normal,
target chart, target detector, physical clock, or target-fit coefficient.
It does find real primitive debt: why $\tau$ exists, why six distinct $X_A$
are selected, whether their equivalence class is source-derived, whether a
reduction can select one robust cone and scale, and whether the operational
regular locus persists under a later coupled law.  Those debts block adoption
and $g_{\mathrm{eff}}$.

\section{Refuter cases and stop conditions}

The following countercases are part of the result, not deferred prose.
\begin{enumerate}[leftmargin=*]
  \item If $q^6=0$, the inherited ratios and token update are undefined.
  \item If $q^2=H(q^1)$ locally, then
        $\dd\rho_1\wedge\dd\rho_2$ may vanish and the operational rank screen
        fails even though the transport equations can still hold.
  \item If $\dd\tau(X_A)=0$, $X_A$ is tangent to the proposed initial surface
        and the stated Cauchy construction fails for that channel.
  \item If an $X_A$ is not locally Lipschitz, uniqueness of characteristics
        can fail; such data are inadmissible.
  \item Prescribing unrelated values on an outflow boundary overdetermines the
        characteristic solution.
  \item Replacing $X_A$ after comparison with a GR target is target fitting and
        invalidates the source-purity claim.
  \item Contracting derivatives with an undeclared inverse tensor or writing
        an action with an undeclared density imports hidden geometry.
  \item A constant solution satisfies the equations but, by itself, fails the
        nontrivial regular-background acceptance burden.
  \item Restricting all six $q^A$ to functions of one amplitude defeats the
        rank-two operational witness and is not an accepted subfamily.
  \item Calling \eqref{eq:balance} a physical conservation law, or
        \eqref{eq:tokenupdate} a realistic matter law, is an invalid overread.
\end{enumerate}

\section{Distance-to-GR matrix and bridge status}

\begin{center}
\small
\begin{tabular}{p{0.29\linewidth}p{0.64\linewidth}}
\toprule
Burden & Status after this packet \\
\midrule
Source ontology primitives & B1 state remains proposal-only; $\tau$ and $X_A$ add explicit primitive debt and are not adopted. \\
Source equivalence $\mathrm{EqSrc}$ & Passive source-chart naturality is exact; no physical equivalence law is adopted. \\
$\mathrm{RetainH}$ & Unchanged and not derived. \\
$\mathrm{GenH}$ & Unchanged and not derived. \\
$\mathrm{ObsLoc}_{\mathrm{lc}}$ & The inherited two-ratio operational patch has a nonconstant dynamical witness. \\
$\mathrm{Resp}_{\mathrm{lc}}$ & Rank two is retained on the explicit background; this remains necessary, not sufficient. \\
$M_{\mathrm{src}}$ & The primitive source arena is still scaffolding debt; no manifold adoption occurs. \\
$g_{\mathrm{eff}}$ & Blocked: the degree-six product of linear factors supplies no unique quadratic cone or scale. \\
Matter coupling & Blocked; the conditional P7 update is not a matter law. \\
Einstein equations & Blocked. \\
Finite-variation robustness & Not established; P3--T03 must define refinement and error control. \\
Benchmark promotion & Blocked. \\
Gate Chair status & No protected verdict requested or issued. \\
Current route freeze & Not frozen: a materially new explicit source model was constructed; next work must test controlled continuum refinement. \\
Hard-fail status & No hidden target import in the accepted candidate; every listed smuggling variant is rejected. \\
\bottomrule
\end{tabular}
\end{center}

The bridge attempt is explicit:
\begin{equation}
  (\D_U,q)\longmapsto
  \sigma_{\D}(\xi)=\operatorname{diag}
  \bigl(\xi(X_1),\ldots,\xi(X_6)\bigr)
  \longmapsto
  \mathcal P_{\D}(\xi)=\prod_A\xi(X_A).
\end{equation}
The first map is constructed and source-native.  The attempted continuation
from $\mathcal P_{\D}$ to one physical cone, scale, and
$g_{\mathrm{eff}}$ is not constructed: it requires a new source-derived
reduction/selection object and later robustness evidence.

\section{Decisive result}

The Candidate Constructor result is \emph{constructed candidate}:
`CAND-V22-B1-SIX-TRANSPORT-DYNAMICS-V1` consists of
\eqref{eq:datum}--\eqref{eq:eom} with the declared admissibility and data
rules.  It has a nonconstant regular solution, exact local characteristic
well-posedness, an exact linearization, and a source-derived degree-six
principal polynomial.  It neither adopts its coefficient datum nor constructs
a unique physical cone, scale, or effective metric.  The next admissible V22
packet after checkpoint is P3--T03, which must define controlled continuum
refinement and error criteria for this proposal-only dynamics without changing
its coefficients by target fitting.

In particular, this packet does not construct an effective metric.

\end{document}
